
\section{Conclusion and Outlook} \label{conclusion}

In conclusion, in this project a complete end-to-end 3D multi-person tracking pipeline was designed, implemented, and evaluated. 
The quantitative evaluation on the CMU Panoptic dataset provided some architectural insight. Spatial viewpoint density exhibits a far higher return on tracking accuracy than increasing deep learning model capacity. Scaling the camera view count from 4 to 16 with the YOLO26n-pose model reduced the mean per-joint position error by more than half, from 5.44 cm down to 2.58 cm. In contrast, replacing this lightweight model with the heavy YOLO26l-pose model under a fixed 4-camera setup yielded smaller accuracy gains at the cost of a massive computational bottleneck, dropping system throughput to approximately 3 FPS.

Despite the relatively high localization accuracy of the system under optimal conditions, there remain several limitations. For example, the process noise matrix, the measurement noise boundaries, the gating thresholds and the physics model were all hand-tuned based on empirical intuition rather than optimized directly from dataset statistics. This restricts the filter's achievable performance and creates a manual optimization bottleneck. Furthermore, the system models the tracking state of each joint mostly independently. Since the system lacks explicit physical constraints, noisy or partially occluded joint detections can lead to anatomically impossible skeletal deformations. Finally, as demonstrated in the qualitative evaluations, tracking quality degrades when targets are partially occluded or not in view of the camera views. Since 3D initialization requires triangulation across at least two views, the system exhibits tracking delays during entry and spatial drifting during single-camera visibility.

Possible future improvements to address these limitations could include the learning of system parameters from the dataset. The parameters could be optimized in an end-to-end differentiable manner using backpropagation through time on a set of training sequences. A further improvement could be the addition of anatomical plausibility constraints on the state predictions. These constraints could possibly be included as pseudo-measurements in the update step of the extended Kalman filter. Another avenue of investigation could be the fine-tuning of the YOLO based 2D pose estimation models. Also, the addition of additional information sources, such as monocular depth values could make it possible to track people even using single camera views. Exploring more advanced filtering methods such as particle filters, potentially also using occlusion aware likelihoods could be promising future directions. Finally, the evaluation in this project has been limited by time constraints. As such, the evaluation of the system using a larger set of more diverse sequences would be highly valuable to uncover further areas for improvement.
